\section{Ví dụ thực tiễn các nước Xã hội chủ nghĩa}

\subsection{Đặt vấn đề}

Khi phân tích quy luật giá trị thặng dư trong các quốc gia xã hội chủ nghĩa hiện nay, cần đặt vấn đề trong bối cảnh của thời kỳ quá độ. Các quốc gia này không còn vận hành theo mô hình kế hoạch hóa tập trung khép kín như trước, mà tồn tại nền kinh tế hàng hóa nhiều thành phần, có thị trường, có doanh nghiệp nhà nước, doanh nghiệp tư nhân và khu vực có vốn đầu tư nước ngoài. Vì vậy, phần giá trị dôi ra ngoài chi phí sức lao động vẫn tiếp tục được tạo ra trong quá trình sản xuất, nhất là ở những khu vực còn tồn tại quan hệ lao động làm thuê. Điểm cần nhấn mạnh không chỉ là có thặng dư hay không, mà là phần thặng dư ấy thuộc về ai, được phân phối như thế nào và phục vụ mục tiêu nào.

Xét từ lập trường Mác -- Lênin, mục tiêu của chủ nghĩa xã hội không phải là duy trì cơ chế để một nhóm tư nhân chiếm đoạt phần giá trị do người lao động tạo ra, mà là từng bước hạn chế và tiến tới xóa bỏ cơ chế chiếm đoạt ấy, đồng thời khẳng định vai trò trung tâm của người lao động trong quá trình phát triển. Vì thế, ở các quốc gia xã hội chủ nghĩa hiện nay, dù cơ chế thị trường vẫn tồn tại, nhà nước vẫn giữ vai trò lớn trong việc nắm giữ các lĩnh vực then chốt và điều tiết một phần đáng kể thặng dư vào đầu tư công, giáo dục, y tế, an sinh xã hội và hạ tầng. Chính ở điểm đó mới thấy rõ sự khác biệt căn bản với chủ nghĩa tư bản, nơi mục tiêu trực tiếp của sản xuất là tối đa hóa lợi nhuận tư nhân \cite{vdtt2}.

\subsection{Trung Quốc}

Trung Quốc là trường hợp tiêu biểu nhất cho mô hình kinh tế thị trường mang định hướng xã hội chủ nghĩa. Theo số liệu của Cục Thống kê Quốc gia Trung Quốc, GDP năm 2024 đạt 134,8066 nghìn tỷ nhân dân tệ, tăng 5,0\% so với năm trước \cite{vdtt1}. Trong cùng năm, tổng lợi nhuận của các doanh nghiệp công nghiệp trên quy mô được thống kê đạt 7.431,05 tỷ nhân dân tệ; trong đó doanh nghiệp nhà nước nắm quyền chi phối đạt 2.139,73 tỷ nhân dân tệ, còn doanh nghiệp tư nhân đạt 2.324,58 tỷ nhân dân tệ \cite{vdtt2}. Những con số này cho thấy khu vực nhà nước và khu vực tư nhân cùng tồn tại, cùng tham gia vào quá trình tích lũy kinh tế.

Điều đáng chú ý là phần thặng dư kinh tế ở Trung Quốc không chỉ dừng lại ở lợi nhuận doanh nghiệp. Năm 2024, chi ngân sách của Trung Quốc đạt 28,46 nghìn tỷ nhân dân tệ; riêng chi cho giáo dục là 4,21 nghìn tỷ nhân dân tệ, còn chi cho an sinh xã hội và việc làm cũng ở mức 4,21 nghìn tỷ nhân dân tệ \cite{vdtt3}. Điều đó cho thấy một phần kết quả kinh tế đã được nhà nước huy động và phân bổ lại cho các mục tiêu xã hội, chứ không hoàn toàn chuyển thành lợi nhuận tư nhân như trong mô hình tư bản chủ nghĩa.

Dưới góc nhìn Mác -- Lênin, quy luật giá trị thặng dư ở Trung Quốc không mất đi, nhưng hình thức chiếm hữu và sử dụng phần thặng dư đã thay đổi. Trong khu vực tư nhân và khu vực có vốn đầu tư nước ngoài, logic lợi nhuận vẫn vận hành khá rõ; song khu vực nhà nước vẫn nắm giữ nhiều lĩnh vực quan trọng, từ đó một phần lớn thặng dư xã hội được chuyển vào đầu tư công và chi tiêu xã hội. Tuy vậy, mâu thuẫn không phải đã biến mất hoàn toàn. Theo World Bank, nếu dùng chuẩn nghèo của nhóm nước thu nhập trung bình cao là 6,85 USD/ngày (PPP 2017), năm 2021 Trung Quốc vẫn còn khoảng 17\% dân số ở dưới ngưỡng này \cite{vdtt4}. Điều đó cho thấy thị trường vẫn có thể tạo ra phân hóa xã hội, nên vai trò điều tiết của nhà nước vẫn giữ ý nghĩa quyết định.

\subsection{Việt Nam}

Việt Nam là trường hợp điển hình của nền kinh tế thị trường định hướng xã hội chủ nghĩa trong thời kỳ quá độ. Theo số liệu công bố đầu năm 2025, GDP năm 2024 tăng 7,09\% \cite{vdtt5}. Cùng năm, tổng vốn FDI đăng ký đạt gần 38,23 tỷ USD, vốn thực hiện đạt khoảng 25,35 tỷ USD; khu vực FDI chiếm gần 71,7\% kim ngạch xuất khẩu của cả nước \cite{vdtt6}. Những số liệu này cho thấy nền kinh tế Việt Nam hội nhập rất sâu vào thị trường thế giới, đặc biệt ở khu vực công nghiệp chế biến, chế tạo và xuất khẩu.

Ở khu vực kinh tế nhà nước, Petrovietnam là ví dụ khá tiêu biểu. Năm 2024, tập đoàn này đạt doanh thu vượt 1 triệu tỷ đồng và nộp ngân sách khoảng 165 nghìn tỷ đồng, tương đương gần 9\% tổng thu ngân sách nhà nước \cite{vdtt7}. Ở bình diện xã hội, báo cáo của Chính phủ cho biết giai đoạn 2021--2025, Việt Nam đã chi khoảng 1,1 triệu tỷ đồng cho an sinh xã hội, tương đương khoảng 17\% tổng chi ngân sách; tỷ lệ hộ nghèo đa chiều giảm từ 4,4\% năm 2021 xuống 1,3\% năm 2025 \cite{vdtt8}. World Bank cũng ghi nhận tỷ lệ dân số sống dưới mức 3,65 USD/ngày đã giảm từ 14\% năm 2010 xuống dưới 4\% năm 2023 \cite{vdtt9}. Những chỉ dấu đó cho thấy một phần quan trọng của kết quả tăng trưởng đã được chuyển hóa thành phúc lợi và điều kiện sống cho số đông dân cư.

Theo lý luận Mác -- Lênin, khu vực tư nhân và khu vực FDI ở Việt Nam vẫn là nơi quy luật giá trị thặng dư biểu hiện khá rõ, vì ở đó quan hệ lao động làm thuê vẫn tồn tại và doanh nghiệp vẫn chịu áp lực lợi nhuận. Tuy nhiên, Việt Nam không phải là nền kinh tế tư bản chủ nghĩa theo nghĩa đầy đủ. Sự khác biệt nằm ở chỗ nhà nước vẫn giữ vai trò chủ đạo trong một số ngành then chốt và dùng ngân sách, doanh nghiệp nhà nước cùng các chính sách an sinh để điều tiết lại kết quả tăng trưởng. Nói cách khác, một bộ phận thặng dư vẫn được tạo ra theo cơ chế thị trường, nhưng một bộ phận quan trọng của nó được nhà nước thu hồi và phân phối lại nhằm phục vụ hạ tầng, phúc lợi và ổn định xã hội \cite{vdtt6}.

\subsection{Lào}

Lào là một quốc gia xã hội chủ nghĩa có quy mô kinh tế nhỏ hơn nhiều so với Trung Quốc và Việt Nam, nhưng vẫn là một trường hợp đáng chú ý. Theo World Bank, kinh tế Lào năm 2024 tăng khoảng 4,1\%, chủ yếu nhờ du lịch, vận tải, điện năng, khai khoáng, nông nghiệp và một phần sản xuất chế biến \cite{vdtt10}. Tuy nhiên, tăng trưởng này diễn ra trong bối cảnh lạm phát còn cao, nợ công ở mức không bền vững và áp lực vĩ mô vẫn rất lớn \cite{vdtt10}.

Dữ liệu thương mại chính thức cho thấy trong tháng 11/2024, tổng kim ngạch xuất khẩu của Lào (không tính điện) đạt khoảng 610 triệu USD; các nhóm hàng xuất khẩu chính gồm vàng, thiết bị điện, cao su và kali \cite{vdtt11}. Nhưng đi cùng với đó, World Bank cho biết lạm phát cao và sức ép kinh tế đã làm suy giảm sức mua của hộ gia đình, buộc nhiều gia đình phải cắt giảm chi tiêu thiết yếu, trong đó có giáo dục và y tế \cite{vdtt12}. Điều này cho thấy phần kết quả kinh tế tạo ra chưa được chuyển hóa đầy đủ thành cải thiện đời sống xã hội.

Dưới góc nhìn Mác -- Lênin, trường hợp Lào cho thấy sự tồn tại của định hướng xã hội chủ nghĩa chưa tự động bảo đảm rằng phần thặng dư kinh tế sẽ được phân phối hiệu quả cho người lao động. Trong điều kiện nền kinh tế còn nhỏ, phụ thuộc nhiều vào xuất khẩu tài nguyên và năng lượng, khả năng điều tiết của nhà nước còn hạn chế. Vì thế, gánh nặng điều chỉnh dễ rơi vào đời sống dân cư. Lào là ví dụ cho thấy khoảng cách giữa mục tiêu xã hội chủ nghĩa và điều kiện vật chất thực tế của thời kỳ quá độ.

\subsection{Cuba}

Cuba là trường hợp đặc biệt trong số các quốc gia xã hội chủ nghĩa hiện nay. Theo ECLAC, GDP của Cuba năm 2024 ước giảm 1,1\%, phản ánh tình trạng tăng trưởng yếu trong bối cảnh thiếu hụt ngoại tệ, khó khăn nhập khẩu và sức ép kinh tế kéo dài \cite{vdtt13}. Khác với Trung Quốc hay Việt Nam, Cuba không nổi bật ở quy mô tích lũy kinh tế, mà chủ yếu thể hiện ở nỗ lực giữ định hướng xã hội trong điều kiện nguồn lực vật chất hạn chế.

Điểm nổi bật của Cuba là ưu tiên ngân sách cho các lĩnh vực xã hội. Theo Prensa Latina, 46\% tổng chi ngân sách được tập trung vào y tế và giáo dục \cite{vdtt14}. UNESCO cũng xác nhận Cuba tiếp tục tham gia chương trình ERCE 2025, cho thấy giáo dục vẫn được duy trì như một ưu tiên chính sách quan trọng \cite{vdtt15}. Trong điều kiện tăng trưởng thấp, việc vẫn ưu tiên mạnh cho y tế và giáo dục cho thấy phần nguồn lực mà nhà nước nắm giữ không được phân bổ hoàn toàn theo logic lợi nhuận.

Từ góc nhìn Mác -- Lênin, Cuba cho thấy một dạng khác của vấn đề: thặng dư kinh tế có thể không lớn, nhưng định hướng phân phối lại vẫn rất rõ. Tuy nhiên, trường hợp này cũng cho thấy một giới hạn quan trọng: nếu lực lượng sản xuất phát triển chậm và quy mô của cải vật chất không đủ lớn, thì việc duy trì phúc lợi phổ quát sẽ ngày càng khó khăn. Điều đó cho thấy định hướng xã hội chủ nghĩa về phân phối không thể tách rời yêu cầu phát triển lực lượng sản xuất \cite{vdtt13}.

\subsection{Đánh giá chung}

Qua bốn trường hợp trên có thể thấy, trong các quốc gia xã hội chủ nghĩa hiện nay, quy luật giá trị thặng dư không nên bị hiểu theo lối giản đơn. Ở Trung Quốc và Việt Nam, nơi thị trường phát triển mạnh và khu vực tư nhân, FDI có vai trò lớn, cơ chế tạo ra giá trị thặng dư vẫn tồn tại khá rõ trong các quan hệ lao động làm thuê. Ở Lào, thặng dư kinh tế có tồn tại nhưng khả năng điều tiết còn hạn chế nên hiệu quả phân phối lại chưa cao. Ở Cuba, quy mô thặng dư vật chất không lớn, nhưng định hướng phân phối xã hội vẫn thể hiện rõ qua ưu tiên dành cho y tế và giáo dục. Sự khác biệt cơ bản với các quốc gia tư bản chủ nghĩa nằm ở chỗ phần thặng dư ấy không hoàn toàn bị tư nhân chiếm hữu, mà còn được nhà nước điều tiết để phục vụ các mục tiêu công cộng \cite{vdtt3}.

\subsection{So sánh với các quốc gia tư bản chủ nghĩa}

So với các quốc gia tư bản chủ nghĩa, các quốc gia xã hội chủ nghĩa hiện nay vẫn có những điểm giống nhau về hình thức vận động kinh tế: đều có sản xuất hàng hóa, tiền tệ, cạnh tranh, giá cả thị trường và ở nhiều khu vực vẫn tồn tại quan hệ lao động làm thuê. Vì vậy, trong khu vực tư nhân và FDI của Trung Quốc hay Việt Nam, quy luật giá trị thặng dư vẫn có thể biểu hiện khá gần với các nền kinh tế tư bản chủ nghĩa \cite[91-93]{gt}.

Tuy nhiên, khác biệt nằm ở chỗ trong chủ nghĩa tư bản, mục tiêu trực tiếp của sản xuất là tối đa hóa lợi nhuận tư nhân; còn trong các quốc gia xã hội chủ nghĩa, nhà nước vẫn tìm cách giữ vai trò trung tâm trong việc phân phối một phần thặng dư vào hạ tầng, an sinh, giáo dục và y tế. Chính điểm đó làm thay đổi tính chất xã hội của quá trình tích lũy \cite{vdtt2}.

\subsection{Kết luận}

Phân tích các quốc gia xã hội chủ nghĩa hiện nay cho thấy quy luật giá trị thặng dư vẫn còn giá trị giải thích, nhưng phải được vận dụng một cách cụ thể. Không thể nói đơn giản rằng trong các nước xã hội chủ nghĩa không còn thặng dư, cũng không thể đồng nhất mọi hình thức thặng dư với bóc lột tư bản chủ nghĩa theo nghĩa kinh điển. Điều quan trọng là phần thặng dư ấy được tạo ra trong cấu trúc sở hữu nào, được nhà nước điều tiết đến đâu, và cuối cùng có được sử dụng để phục vụ lợi ích chung của xã hội hay không. Trung Quốc, Việt Nam, Lào và Cuba cho thấy bốn cách biểu hiện khác nhau của cùng một vấn đề, đồng thời cho thấy rằng khi nghiên cứu các quốc gia xã hội chủ nghĩa hiện nay, cần nhìn vào mối quan hệ giữa thị trường, sở hữu, nhà nước và phân phối xã hội, thay vì áp dụng máy móc khuôn mẫu của chủ nghĩa tư bản cổ điển \cite{vdtt2}.
